% ============================================================
\chapter{Convergence Theorems}
\label{ch:convergence}
% ============================================================

\section{Introduction}

Here is the question every analyst eventually faces: when can one interchange a limit and an integral? The Riemann integral offers a spartan answer: under uniform convergence, and that is all. But in practice --- Fourier series, probability theory, statistical mechanics --- uniform convergence is a rare luxury. This is precisely where the Lebesgue integral reveals its full power. Between 1902 and 1910, Lebesgue, then Beppo Levi and Pierre Fatou, established three convergence theorems that form the beating heart of modern integration. Beppo Levi's monotone convergence theorem (1906) handles increasing sequences; Fatou's lemma (1906) provides a universal inequality; and Lebesgue's dominated convergence theorem crowns the edifice by allowing the limit-integral exchange whenever an integrable function ``dominates'' the sequence. These three results are, without exaggeration, the most frequently used tools in all of modern analysis.

\section{Monotone convergence theorem (Beppo Levi)}

\begin{theorem}[Monotone convergence]
\label{thm:monotone_convergence}
Let $(f_n)_{n \in \N}$ be an \emph{increasing} sequence of non-negative measurable
functions: $0 \leq f_1 \leq f_2 \leq \ldots$ Let $f = \lim_n f_n = \sup_n f_n$.
Then:
\[
\int_X f\, d\mu = \lim_{n \to \infty} \int_X f_n\, d\mu
= \sup_{n \in \N} \int_X f_n\, d\mu.
\]
\end{theorem}

\begin{proof}
Set $I = \lim_n \int f_n\, d\mu$ (the limit exists in $[0, +\infty]$ since the
sequence $(\int f_n\, d\mu)$ is increasing by monotonicity of the integral).

\emph{Inequality $\int f\, d\mu \geq I$}: since $f_n \leq f$ for all $n$,
$\int f_n\, d\mu \leq \int f\, d\mu$, so $I \leq \int f\, d\mu$.

\emph{Inequality $\int f\, d\mu \leq I$}: let $\varphi$ be a measurable simple
function with $0 \leq \varphi \leq f$ and $\alpha \in (0, 1)$. Set
$E_n = \{x : f_n(x) \geq \alpha \varphi(x)\}$. Then $E_n \uparrow X$ (since for
every $x$, $f_n(x) \uparrow f(x) \geq \varphi(x) > \alpha \varphi(x)$ eventually).
We have:
\[
\int f_n\, d\mu \geq \int_{E_n} f_n\, d\mu \geq \alpha \int_{E_n} \varphi\, d\mu.
\]
Letting $n \to \infty$, by monotone continuity of the measure
$A \mapsto \int_A \varphi\, d\mu$:
\[
I \geq \alpha \int_X \varphi\, d\mu.
\]
Letting $\alpha \to 1^-$: $I \geq \int_X \varphi\, d\mu$. Taking the supremum over
all simple functions $\varphi \leq f$: $I \geq \int f\, d\mu$.
\end{proof}

\begin{keyformulas}[title=\textbf{Monotone convergence}]
If $0 \leq f_n \uparrow f$, then
$\displaystyle\int f_n\, d\mu \uparrow \int f\, d\mu$.
\end{keyformulas}

\begin{corollary}[Series of non-negative functions]
Let $(g_n)_{n \in \N}$ be a sequence of non-negative measurable functions. Then:
\[
\int_X \sum_{n=0}^{\infty} g_n\, d\mu = \sum_{n=0}^{\infty} \int_X g_n\, d\mu.
\]
\end{corollary}

\begin{proof}
Apply the monotone convergence theorem to $f_N = \sum_{n=0}^N g_n$.
\end{proof}

\begin{corollary}[Additivity of the integral]
For measurable $f, g : X \to [0, +\infty]$:
$\int (f + g)\, d\mu = \int f\, d\mu + \int g\, d\mu$.
\end{corollary}

\begin{proof}
Let $(\varphi_n)$ and $(\psi_n)$ be increasing sequences of non-negative simple
functions with $\varphi_n \uparrow f$ and $\psi_n \uparrow g$. Then
$\varphi_n + \psi_n \uparrow f + g$ and
$\int (\varphi_n + \psi_n)\, d\mu = \int \varphi_n\, d\mu + \int \psi_n\, d\mu$.
Apply the monotone convergence theorem to both sides.
\end{proof}

\section{Fatou's lemma}

\begin{theorem}[Fatou's lemma]
\label{thm:fatou}
Let $(f_n)_{n \in \N}$ be a sequence of non-negative measurable functions. Then:
\[
\int_X \liminf_{n \to \infty} f_n\, d\mu \leq
\liminf_{n \to \infty} \int_X f_n\, d\mu.
\]
\end{theorem}

\begin{proof}
Set $g_n = \inf_{k \geq n} f_k$. Then $g_n$ is non-negative measurable,
$g_n \leq f_n$, and $g_n \uparrow \liminf_n f_n$. By the monotone convergence
theorem:
\[
\int \liminf_n f_n\, d\mu = \lim_n \int g_n\, d\mu \leq
\liminf_n \int f_n\, d\mu,
\]
the last inequality following from $\int g_n\, d\mu \leq \int f_n\, d\mu$.
\end{proof}

\begin{warning}[title=\textbf{The inequality can be strict}]
Take $f_n = n \cdot \mathbf{1}_{(0, 1/n)}$ on $(\R, \lambda)$. Then $f_n \to 0$
a.e. but $\int f_n\, d\lambda = 1$ for all $n$. So
$0 = \int \liminf f_n\, d\lambda < \liminf \int f_n\, d\lambda = 1$.
\end{warning}

\begin{corollary}[Reverse Fatou lemma]
If $(f_n)$ is a sequence of non-negative measurable functions and there exists
$g \in L^1(\mu)$ with $f_n \leq g$ a.e. for all $n$, then:
\[
\limsup_{n \to \infty} \int_X f_n\, d\mu \leq
\int_X \limsup_{n \to \infty} f_n\, d\mu.
\]
\end{corollary}

\begin{proof}
Apply Fatou's lemma to the sequence $g - f_n \geq 0$.
\end{proof}

\section{Dominated convergence theorem (Lebesgue)}

\begin{theorem}[Dominated convergence]
\label{thm:dominated_convergence}
Let $(f_n)$ be a sequence of measurable functions $X \to \overline{\R}$ such that:
\begin{enumerate}
    \item[(i)] $f_n \to f$ $\mu$-a.e. (pointwise convergence almost everywhere);
    \item[(ii)] there exists $g \in L^1(\mu)$, $g \geq 0$, with
    $\abs{f_n} \leq g$ $\mu$-a.e. for all $n$ (\emph{domination}).
\end{enumerate}
Then $f \in L^1(\mu)$, each $f_n \in L^1(\mu)$, and:
\[
\lim_{n \to \infty} \int_X f_n\, d\mu = \int_X f\, d\mu.
\]
Moreover, $\int_X \abs{f_n - f}\, d\mu \to 0$.
\end{theorem}

\begin{proof}
Since $\abs{f_n} \leq g$ a.e. and $f_n \to f$ a.e., we have $\abs{f} \leq g$ a.e.,
so $f \in L^1(\mu)$.

Consider the functions $g + f_n \geq 0$ and $g - f_n \geq 0$. By Fatou's lemma:
\begin{align*}
\int (g + f)\, d\mu &\leq \liminf_n \int (g + f_n)\, d\mu
= \int g\, d\mu + \liminf_n \int f_n\, d\mu, \\
\int (g - f)\, d\mu &\leq \liminf_n \int (g - f_n)\, d\mu
= \int g\, d\mu - \limsup_n \int f_n\, d\mu.
\end{align*}
Since $\int g\, d\mu < \infty$, we deduce:
\[
\int f\, d\mu \leq \liminf_n \int f_n\, d\mu \leq
\limsup_n \int f_n\, d\mu \leq \int f\, d\mu.
\]

For the last assertion, apply dominated convergence to $\abs{f_n - f} \leq 2g$
and $\abs{f_n - f} \to 0$ a.e.
\end{proof}

\begin{keyformulas}[title=\textbf{Dominated convergence theorem}]
\[
\boxed{f_n \to f \text{ a.e.}, \quad \abs{f_n} \leq g \in L^1
\implies \int f_n\, d\mu \to \int f\, d\mu.}
\]
\end{keyformulas}

\section{Applications}

\begin{corollary}[Continuity under the integral sign]
\label{cor:continuity_integral}
Let $f : X \times I \to \R$ where $I \subset \R$ is an interval. Suppose:
\begin{enumerate}
    \item for every $t \in I$, $x \mapsto f(x, t)$ is measurable;
    \item for $\mu$-a.e. $x$, $t \mapsto f(x, t)$ is continuous at $t_0$;
    \item there exists $g \in L^1(\mu)$ with $\abs{f(x, t)} \leq g(x)$ for all
    $t \in I$ and $\mu$-a.e. $x$.
\end{enumerate}
Then $F(t) = \int_X f(x, t)\, d\mu(x)$ is continuous at $t_0$.
\end{corollary}

\begin{proof}
If $t_n \to t_0$, then $f(x, t_n) \to f(x, t_0)$ a.e. with $\abs{f(x, t_n)}
\leq g(x)$. By dominated convergence, $F(t_n) \to F(t_0)$.
\end{proof}

\begin{corollary}[Differentiation under the integral sign]
\label{cor:differentiation_integral}
Let $f : X \times I \to \R$ where $I$ is an open interval. Suppose:
\begin{enumerate}
    \item for every $t \in I$, $x \mapsto f(x, t)$ is integrable;
    \item for $\mu$-a.e. $x$, $t \mapsto f(x, t)$ is differentiable on $I$;
    \item there exists $g \in L^1(\mu)$ with
    $\abs{\frac{\partial f}{\partial t}(x, t)} \leq g(x)$ for all $t \in I$
    and $\mu$-a.e. $x$.
\end{enumerate}
Then $F(t) = \int_X f(x, t)\, d\mu(x)$ is differentiable on $I$ and:
\[
F'(t) = \int_X \frac{\partial f}{\partial t}(x, t)\, d\mu(x).
\]
\end{corollary}

\begin{proof}
For $h \neq 0$, $\frac{F(t+h) - F(t)}{h} = \int \frac{f(x, t+h) - f(x, t)}{h}
\, d\mu(x)$. By the mean value theorem,
$\abs{\frac{f(x, t+h) - f(x, t)}{h}} \leq g(x)$. Letting $h \to 0$, dominated
convergence gives the result.
\end{proof}

\section{Convergence in $L^1$}

\begin{theorem}[Characterization of $L^1$ convergence]
Let $f_n, f \in L^1(\mu)$. The following are equivalent:
\begin{enumerate}
    \item $\int \abs{f_n - f}\, d\mu \to 0$ ($L^1$ convergence).
    \item $f_n \to f$ in measure and the family $(f_n)$ is uniformly integrable.
\end{enumerate}
\end{theorem}

\begin{definition}[Uniform integrability]
A family $(f_i)_{i \in I}$ of integrable functions is \emph{uniformly integrable}
if:
\[
\lim_{M \to \infty} \sup_{i \in I} \int_{\{\abs{f_i} > M\}} \abs{f_i}\, d\mu = 0.
\]
\end{definition}

\begin{theorem}[Vitali]
\label{thm:vitali_convergence}
Let $\mu(X) < \infty$. If $(f_n)$ is uniformly integrable and $f_n \to f$ in
measure, then $f \in L^1(\mu)$ and $\int \abs{f_n - f}\, d\mu \to 0$.
\end{theorem}

\section{Exercises}

\begin{exercise}
Compute $\lim_{n \to \infty} \int_0^{\infty}
\frac{n \sin(x/n)}{x(1 + x^2)}\, dx$, justifying the interchange of limit and
integral.
\end{exercise}

\begin{exercise}
Show that $F(\alpha) = \int_0^{\infty} \frac{e^{-\alpha x}}{1 + x^2}\, dx$ is
$C^\infty$ on $(0, \infty)$ and compute $F'(\alpha)$.
\end{exercise}

\begin{exercise}
Let $f \in L^1(\R)$. Show that $F(t) = \int_{\R} e^{itx} f(x)\, dx$ is continuous
on $\R$ (this is the Fourier transform of $f$).
\end{exercise}

\begin{exercise}
Give an example of a sequence $(f_n)$ in $L^1([0,1], \lambda)$ with $f_n \to 0$
a.e. but $\int f_n\, d\lambda \not\to 0$. Why does the dominated convergence
theorem not apply?
\end{exercise}

\begin{exercise}
Compute $\lim_{n \to \infty} \int_0^1 (1 + x/n)^n e^{-2x}\, dx$.
\end{exercise}

\begin{exercise}
Let $f \in L^1(\R)$. Show that
$\lim_{h \to 0} \int_{\R} \abs{f(x + h) - f(x)}\, dx = 0$
(continuity of translation in $L^1$).
\end{exercise}
